\renewcommand{\thesection}{S\arabic{section}}
\setcounter{section}{1}

\subsection{Discrete one-sided evaluation of the separability functional}
\label{sec:Appendix_S2}

The separability functional in Eq.~\eqref{eq:J-def} is defined using a
two-sided continuous Fourier representation. In numerical calculations,
however, the upstream trajectories and delay kernels are real-valued,
discretely sampled sequences. Their discrete Fourier transforms therefore
satisfy conjugate symmetry, allowing the equivalent calculation to be
performed using only the nonnegative-frequency coefficients.

Let
\[
\Delta U_n
=
U_{1,n}-U_{0,n},
\qquad
n=0,\ldots,n_U-1,
\]
denote the discrete upstream contrast, and let
\[
g_n,
\qquad
n=0,\ldots,n_g-1,
\]
denote the discrete delay probability mass function. The corresponding
difference between the expected downstream trajectories is the linear
convolution
\[
\Delta D_n
=
p\,(\Delta U*g)_n.
\]

To represent this linear convolution using the discrete Fourier transform,
both sequences are zero-padded to the full convolution length
\begin{equation}
L=n_U+n_g-1.
\label{eq:full-convolution-length}
\end{equation}
This padding prevents the circular wraparound that would otherwise arise
from multiplication in the discrete Fourier domain.

For a sampling interval \(\Delta t\), define the length-\(L\) discrete
Fourier transforms
\begin{align}
\Delta\widetilde U_m
&=
\sum_{n=0}^{L-1}
\Delta U_n
\exp\left(
-\frac{2\pi i mn}{L}
\right),
\\
\widetilde G_m
&=
\sum_{n=0}^{L-1}
g_n
\exp\left(
-\frac{2\pi i mn}{L}
\right),
\end{align}
evaluated at frequencies
\begin{equation}
f_m
=
\frac{m}{L\Delta t}.
\label{eq:discrete-frequencies}
\end{equation}

Under this convention, the Fourier transform of the downstream contrast is
\begin{equation}
\Delta\widetilde D_m
=
p\,\widetilde G_m\,
\Delta\widetilde U_m.
\label{eq:discrete-downstream-contrast}
\end{equation}

Because \(\Delta U_n\) and \(g_n\) are real-valued,
\begin{equation}
\Delta\widetilde U_{L-m}
=
\Delta\widetilde U_m^{\,*},
\qquad
\widetilde G_{L-m}
=
\widetilde G_m^{\,*},
\end{equation}
where \(^*\) denotes complex conjugation. Consequently, the spectral
integrand is symmetric between positive and negative frequencies. The
two-sided spectral sum can therefore be evaluated from the one-sided real
Fourier transform over
\[
m=0,\ldots,\left\lfloor\frac{L}{2}\right\rfloor.
\]

The discrete one-sided approximation to the model-specific separability
functional is
\begin{equation}
J_{\mathcal M}^{(L)}(U_0,U_1)
=
\frac{\Delta t}{L}
\sum_{m=0}^{\lfloor L/2\rfloor}
q_m
\frac{
\left|p\,\widetilde G_m\right|^2
}{
N_{\mathcal M}(f_m)
}
\left|
\Delta\widetilde U_m
\right|^2,
\label{eq:discrete-one-sided-J}
\end{equation}
where the one-sided multiplicity weights are
\begin{equation}
q_m
=
\begin{cases}
1,
&
m=0,
\\[4pt]
1,
&
L\ \text{is even and}\ m=L/2,
\\[4pt]
2,
&
\text{otherwise}.
\end{cases}
\label{eq:one-sided-weights}
\end{equation}

The zero-frequency coefficient receives unit weight because it has no
distinct negative-frequency counterpart. When \(L\) is even, the Nyquist
coefficient \(m=L/2\) is also self-conjugate and therefore receives unit
weight. Every other positive-frequency coefficient represents a conjugate
pair and consequently receives weight two.

The normalization in Eq.~\eqref{eq:discrete-one-sided-J} follows from
Parseval's identity under the discrete Fourier-transform convention above:
\begin{equation}
\Delta t
\sum_{n=0}^{L-1}
|x_n|^2
=
\frac{\Delta t}{L}
\sum_{m=0}^{L-1}
|\widetilde X_m|^2.
\label{eq:discrete-parseval}
\end{equation}
Thus, the one-sided weighted sum in
Eq.~\eqref{eq:discrete-one-sided-J} recovers exactly the same spectral
energy as the corresponding two-sided discrete sum.

For daily observations, \(\Delta t=1\) day, and
Eq.~\eqref{eq:discrete-one-sided-J} simplifies to
\begin{equation}
J_{\mathcal M}^{(L)}(U_0,U_1)
=
\frac{1}{L}
\sum_{m=0}^{\lfloor L/2\rfloor}
q_m
\frac{
\left|p\,\widetilde G_m\right|^2
}{
N_{\mathcal M}(f_m)
}
\left|
\Delta\widetilde U_m
\right|^2.
\label{eq:discrete-one-sided-J-daily}
\end{equation}

For additive Gaussian noise with a constant standard deviation
\(\sigma_\varepsilon\), the effective noise scale is
\[
N_{\mathcal M}(f_m)=\sigma_\varepsilon^2,
\]
giving
\begin{equation}
J_{\mathrm{G}}^{(L)}(U_0,U_1)
=
\frac{1}{L}
\sum_{m=0}^{\lfloor L/2\rfloor}
q_m
\frac{
\left|p\,\widetilde G_m\right|^2
\left|\Delta\widetilde U_m\right|^2
}{
\sigma_\varepsilon^2
}.
\label{eq:discrete-gaussian-J}
\end{equation}

More generally, a frequency-dependent effective noise scale
\(N_{\mathcal M}(f_m)\) can be inserted directly into
Eq.~\eqref{eq:discrete-one-sided-J}. For the count observation models
considered in this study, the corresponding conservative likelihood-based
variance normalization is constant across frequency, as specified in
Table~\ref{tab:noise-models} and derived in
Appendix~\ref{sec:Appendix_S1}.

Equation~\eqref{eq:discrete-one-sided-J} is therefore the discrete numerical
counterpart of the continuous two-sided functional in
Eq.~\eqref{eq:J-def}. The use of a one-sided real Fourier transform changes
only the computational representation and does not change the value or
interpretation of the separability functional.